\documentclass[10pt,a4paper]{amsart}
\usepackage[margin=19mm]{geometry}
\usepackage{amsmath,amssymb,mathtools}
\usepackage[scr=rsfs,cal=euler]{mathalfa}
\usepackage{xfrac}
\usepackage[unicode,hidelinks,bookmarks=false]{hyperref}
\newcommand{\ie}{i.e.\ }
\newcommand{\cf}{cf.\ }
\newcommand{\resp}{respectively\ }
\newcommand{\Subsection}{\subsection}
\providecommand{\nobreakdash}{\nobreak\hbox{-}}
\setlength{\emergencystretch}{2em}
\multlinegap0pt
\usepackage{bm}
\usepackage{bbm}
\usepackage{dsfont}
\usepackage{xspace}
\usepackage{cancel}
\usepackage{mathtools}
\usepackage{cleveref}
\usepackage{amsthm}
\makeatletter
\@ifundefined{claimproof}{\newtheorem{claimproof}{Claimproof}}{}
\@ifundefined{lem}{\newtheorem{lem}{Lem}}{}
\@ifundefined{thm}{\newtheorem{thm}{Thm}}{}
\@ifundefined{claim}{\newtheorem{claim}[thm]{$\vartriangleright$ Claim}}{}
\makeatother
\makeatletter
\newcommand{\Cc}{\mathscr{C}}
\newcommand{\Dd}{\mathscr{D}}
\newcommand{\Ss}{\mathscr{S}}
\expandafter\ifx\csname claimproof\endcsname\relax
\newenvironment{claimproof}{\begin{proof}[Proof of the claim]
	\renewcommand\qedsymbol{$\vartriangleleft$}
}{\end{proof}}
\fi
\makeatother
\title[On first-order transductions of classes of graphs]{On first-order transductions of classes of graphs}
\author{Samuel Braunfeld, Jaroslav Ne\v{s}et\v{r}il, Patrice Ossona de Mendez, Sebastian Siebertz}
\date{Source: arXiv:2208.14412v8 (CC BY 4.0)}
\begin{document}
\maketitle

\noindent\emph{Source and licence.} On first-order transductions of classes of graphs. Version arXiv:2208.14412v8 (CC BY 4.0), distributed under the Creative Commons Attribution 4.0 International licence, as recorded in the arXiv licence metadata for this exact version and shown on the abstract page https://arxiv.org/abs/2208.14412v8. Full rights notice: licenses/arxiv-2208.14412-CC-BY-4.0.txt.\par\smallskip

\noindent\textit{Editorial scope.} Counted unit: the Thm whose source label is \texttt{thm:rel-normal} in the article (source0.tex lines 849--855, source label \texttt{thm:rel-normal}), delivered together with the article's own complete original proof of it (source0.tex lines 856--930, one \texttt{proof} environment of 75 lines). No mathematics is written, repaired or re-derived: statement and proof are the article's own text line for line. Required context retained uncounted: lem:perturb (lines 753--756); thm:normal (lines 776--782). Every definition, display and label of those lines is retained unchanged. Omitted from this package and not redistributed: 1--752, 757--775, 783--848, 931--2664. Nothing inside a retained excerpt is omitted or altered: every retained body is delivered exactly as the article's lines are written, so no omission receipt of any kind accompanies this package. Citations: the retained excerpts carry no citation command at all, so no bibliography entry of the article is transcribed and no reference is redistributed. Reference adaptation: none was needed and none was made; every reference inside the delivered unit resolves inside it, so no unresolved reference or citation is produced. Printed slips kept literal: no printed word, symbol, hypothesis or number of the counted statement or of its proof was altered. Layout and class: the article's own class is \texttt{lmcs} with packages including lastpage, amssymb, hyperref, cleveref, inputenc, fontenc, microtype, mathrsfs, mathtools, xy, nowidow, relsize, complexity, graphicx; the wrapper is the lane's pinned \texttt{amsart} editorial wrapper at 10pt on A4, which restates the theorem-like environments the retained text needs and adds the article's own macro definitions that the retained text needs (\texttt{\textbackslash Cc}, \texttt{\textbackslash Dd}, \texttt{\textbackslash Ss}), with extra wrapper packages (bm, bbm, dsfont, xspace, cancel, mathtools, cleveref). The delivered numbering is the unit's own. Assets: no retained span contains a figure environment, a graphics-inclusion command, a PDF-page inclusion, a table or a TikZ picture; no image file is copied into this package, the package declares no asset dependency and no image file is redistributed with it. Licence: version arXiv:2208.14412v8 of this article is distributed under the Creative Commons Attribution 4.0 International licence, as recorded in the arXiv licence metadata for this exact version and shown on the abstract page https://arxiv.org/abs/2208.14412v8. A full rights notice is delivered at licenses/arxiv-2208.14412-CC-BY-4.0.txt. Characters: every retained excerpt is pure ASCII, so the delivered engine, XeLaTeX with its default Latin Modern text font, reports no missing character.
\begin{lem}
	\label{lem:perturb}
	For every perturbation $\mathsf P=\oplus Z_1\oplus\dots\oplus Z_k$ and for every graph $G$, $H\in{\mathsf P}(G)$ if and only if there exists a partition $(V_\mathbf x)_{\mathbf x\in \mathbb{F}_2^k}$ of $V(G)$ such that $H$ is obtained from $G$ by flipping the edges $uv$ where $u\in V_{\mathbf x}, v\in V_{\mathbf y}$, and $\langle\mathbf x,\mathbf y\rangle=1$.
\end{lem}

\begin{thm}
	\label{thm:normal}
	Every non-copying transduction $\mathsf T$ is subsumed by the composition of an immersive transduction~$\mathsf T_{\rm imm}$ and a perturbation~$\mathsf P$, that is, $\mathsf T\leq \mathsf P\circ\mathsf T_{\rm imm}$.
	
	Consequently, every transduction $\mathsf T$ is subsumed by the composition of a
	copying operation~$\mathsf C$, an immersive transduction $\mathsf T_{\rm imm}$ and a perturbation~$\mathsf P$, that is, $\mathsf T\leq \mathsf P\circ\mathsf T_{\rm imm}\circ\mathsf C$.
\end{thm}

\begin{thm}
	\label{thm:rel-normal}
	Let  $\mathsf T$ be a non-copying transduction and let a  weakly sparse  class $\Dd$ be a $\mathsf T$-transduction  of a class $\Cc$.
	
	Then, there exist an immersive transduction $\mathsf T_{\rm imm}$, a perturbation $\mathsf P$, and a weakly sparse (hereditary)  class $\Ss$ such that 
	$\mathsf P\circ\mathsf T_{\rm imm}$ subsumes $\mathsf T$ on the chain $(\Cc,\Ss,\Dd)$.
\end{thm}

\begin{proof}
	As the class $\Dd$ is weakly sparse, there exists an integer $t$ such that $K_{t,t}$ is not a subgraph of any graph in $\Dd$.
	
	By \Cref{thm:normal}, the transduction $\mathsf T$ is subsumed by the composition of an immersive transduction $\mathsf T_0$ and a perturbation $\mathsf P_0$.
	Thus, $\Dd\subseteq \mathsf P_0\circ\mathsf T_0(\Cc)$, which rewrites as the property that for every $H\in\Dd$ there exists $G_H\in\Cc$ and $K_H\in\mathsf T_0(G_H)$ such that
	$H\in \mathsf P_0(K_H)$. 
	The problem is that the class $\mathscr K:=\{K_H\colon H\in\Dd\}$ needs not to be weakly sparse.
	
	However, we shall prove that there exists a perturbation $\mathsf P_1$ and an immersive transduction~$\mathsf T_1$ such that, for every $H\in\Dd$, we can choose 
	$K_H'\in\mathsf P_1(K_H)\cap \mathsf T_1(K_H)$ in such a way that the class 
	$\Ss:=\{K_H'\colon H\in\Dd\}$ is weakly sparse.
	Note that $K_H'\in\mathsf P_1(K_H)$ is equivalent to $K_H\in\mathsf P_1(K_H')$, as every perturbation is undone by carrying it out again.
	Then, if we let $\mathsf P=\mathsf P_0\circ\mathsf P_1$ and $\mathsf T_{\rm imm}=\mathsf T_1\circ\mathsf T_0$, the transduction $\mathsf P$ is a perturbation, 
	the transduction $\mathsf T_{\rm imm}$ is an immersive transduction, and, for every $H\in\Dd$, we have
	$K_H'\in \mathsf T_1(K_H)\subseteq \mathsf T_1\circ\mathsf T_0(G_H)=\mathsf T_{\rm imm}(G_H)$ and
	$H\in \mathsf P_0(K_H)\subseteq \mathsf P_0\circ\mathsf P_1(K_H')=\mathsf P(K_H')$.
	Hence, $\mathsf P\circ\mathsf T_{\rm imm}$ subsumes $\mathsf T$ on the chain $(\Cc,\Ss,\Dd)$.
	
	Toward this end, consider  $H\in\Dd$ and let $G_H$ and $K_H$ be such that $G_H\in\Cc$ and 
	$K_H\in\mathsf{T}_0(G_H)\cap\mathsf P_0(H)$.
	Note that this last condition is equivalent to the conditions  
	$K_H\in\mathsf T_0(G_H)$ and $H\in\mathsf P_0(K_H)$.
	
	We analyze how the perturbation $\mathsf P_0$ transforms $K_H$ into $H$.
	According to \Cref{lem:perturb}, there exists a partition $(V_{\mathbf x})_{\mathbf x\in\mathbb F_2^k}$, such that the edges $uv$ of $K_H$ that are flipped to get $H$ are those where $u\in V_{\mathbf x}$, $v\in V_{\mathbf y}$, and $\langle\mathbf{x},\mathbf{y}\rangle=1$.
	%
	For $\mathbf x,\mathbf y\in\mathbb{F}_2^k$, we denote by $B_{\mathbf x,\mathbf y}$ the graph with vertex set $V_{\mathbf x}\cup V_{\mathbf y}$, whose edges are the pairs $uv$ with $u\neq v$, $u\in V_{\mathbf x}$, $v\in V_{\mathbf y}$ and $uv\in E(K_H)$. We say that the pair $(\mathbf x,\mathbf y)$ is \emph{problematic} if $B_{\mathbf x,\mathbf y}$ contains $K_{t,t}$ as a subgraph. Note that if $(\mathbf x,\mathbf y)$ is problematic, then the edges of $B_{\mathbf x,\mathbf y}$  are flipped when transforming $K_H$ into $H$, as this latter graph excludes $K_{t,t}$ as a subgraph. Thus $\langle\mathbf{x},\mathbf{y}\rangle=1$ if $(\mathbf x,\mathbf y)$ is problematic.
	The next claim gives some basic structural properties of problematic pairs.
	
	\begin{claim}
		\label{claim:cc}
		If $(\mathbf x,\mathbf y)$ is problematic, then 
		$B_{\mathbf x,\mathbf y}$ has $c(\mathbf x,\mathbf y)\leq 2t$ connected components, each of diameter at most $6t$.
	\end{claim}
	\begin{claimproof}
		We first consider the case where $\mathbf y=\mathbf x$.
		If $B_{\mathbf x,\mathbf x}$ has $2t$ connected components; then by picking one vertex in each connected component, we get an independent set of size at least $2t$.
		Also, if some connected component of $B_{\mathbf x,\mathbf x}$ has diameter at least $4t$, then $B_{\mathbf x,\mathbf x}$ contains an induced path with $4t$ vertices hence (by picking one vertex over two in this path) an independent set of size $2t$. In both cases, after flipping the edges in $V_{\mathbf x}$ we get a clique of size $2t$, contradicting the assumption that $H$ does not contain $K_{t,t}$ as a subgraph. 
		
		Now consider the case where $\mathbf y\neq\mathbf x$.
		As above, no connected component of $(\mathbf x,\mathbf y)$ has diameter greater than $6t$, as if $v_0,\dots,v_{6t-3}$ is an induced path
		of $B_{\mathbf x,\mathbf y}$, then $\{v_0,v_6,\dots,v_{6t-6}\}$ and 
		$\{v_3,v_9,\dots,v_{6t-3}\}$ define, after flipping, a $K_{t,t}$ subgraph in $H$.	As $(\mathbf x,\mathbf y)$ is problematic, one connected  component $C_0$ of $B_{\mathbf x,\mathbf y}$ contains a $K_{t,t}$ subgraph. Assume for contradiction that $B_{\mathbf x,\mathbf y}$  has $2t-1$ other connected components. Then at least $t$ of them contain a vertex in a same part, say $V_{\mathbf x}$. The edges flipped between these vertices and the vertices in $C_0\cap V_{\mathbf y}$ then contain $K_{t,t}$ as a subgraph, contradicting our assumption.
		\end{claimproof}
	
	We number the connected components of $B_{\mathbf x,\mathbf y}$ from $1$ to $c(\mathbf x,\mathbf y)$ in an arbitrary way (but consistently, when  considering that $B_{\mathbf x,\mathbf y}$ and  $B_{\mathbf y,\mathbf x}$ are the same graph).
	
	We now refine the partition $(V_{\mathbf x})_{\mathbf x\in\mathbb F_2^k}$,
	by partitioning each $V_{\mathbf x}$ into at most $(2t)^{2^k}$ parts~$W_{\mathbf x,f}$, where $f$ maps each vector $\mathbf y$ such that $(\mathbf x,\mathbf y)$ is problematic to $c(\mathbf x,\mathbf y)$ and each vector $\mathbf y$ such that $(\mathbf x,\mathbf y)$ is not problematic to $0$.  
	Note that, for every pair  $(\mathbf x,\mathbf y)$ of vectors, if  
	$W_{\mathbf x,f}$ is any subpart of $V_{\mathbf x}$ and 
	$W_{\mathbf y,g}$ is any subpart of $V_{\mathbf y}$, we have
	\[
	f(\mathbf y)=0\quad\iff\quad(\mathbf x,\mathbf y)\text{ is not problematic}\quad\iff\quad g(\mathbf x)=0.\]
	
	We define a perturbation $\mathsf P_1$ that flips the edges with one end in 
	$W_{\mathbf x,f}$ and one end in~$W_{\mathbf y,g}$ if
	$\langle\mathbf{x},\mathbf{y}\rangle=1$ and
	$f(\mathbf y)\neq 0$ and $f(\mathbf y)=g(\mathbf x)$ (meaning that 	 $(\mathbf x,\mathbf y)$ is a problematic pair, but 	$W_{\mathbf x,f}$ and $W_{\mathbf y,g}$  are in the same connected component of $B_{\mathbf x,\mathbf y}$).
	
	Now, let $K_H'\in\mathsf P_1(K_H)$ be the graph obtained by applying the perturbation with the marking defined above.
	Let $B'_{(\mathbf x,f),(\mathbf y,g)}$ be the subgraph of $K_H'$
	with vertex set  $W_{\mathbf x,f}\cup W_{\mathbf y,g}$ whose edges are those edges of $K_H'$ with one endpoint in  $W_{\mathbf x,f}$ and the other endpoint in~$W_{\mathbf y,g}$.
	\begin{itemize}
		\item If $\langle\mathbf x,\mathbf y\rangle=0$, then $B'_{(\mathbf x,f),(\mathbf y,g)}$ is a subgraph of $H$, thus it does not contain $K_{t,t}$ as a subgraph;
		\item otherwise, if $f(\mathbf y)=0$, then $B'_{(\mathbf x,f),(\mathbf y,g)}$ is a subgraph of $B_{\mathbf x,\mathbf y}$, thus it does not contain $K_{t,t}$ as a subgraph since $f(\mathbf y)=0$ implies that $(\mathbf x,\mathbf y)$ is not problematic;
		\item otherwise, if $f(\mathbf y)=g(\mathbf y)$, then $B'_{(\mathbf x,f),(\mathbf y,g)}$ is a subgraph of $H$, thus it does not contain $K_{t,t}$ as a subgraph;
		\item otherwise,  $B'_{(\mathbf x,f),(\mathbf y,g)}$ is edgeless, as $W_{\mathbf x,f}$ and $W_{\mathbf y,g}$ are included in two different connected components of $B_{\mathbf x,\mathbf y}$.
	\end{itemize}
	As the number of parts is at most $(2^t)^{2^k}$, we deduce that
	$K_H'$ contains no $K_{t',t'}$ as a subgraph, where
	$t'=(2^t)^{2^k}t$.
	
	We now define the immersive transduction $\mathsf T_1$ from the perturbation $\mathsf P_1$ by conditioning the flip of an edge $uv$ to the further condition that the distance between $u$ and $v$ is at most~$6t$. Because the maximum diameter of a connected component of $B_{\mathbf x,\mathbf y}$ is at most $6t$ (by \Cref{claim:cc}), we get $K_H'\in\mathsf T_1(K_H)$ as desired.
\end{proof}

\end{document}
